\documentclass[12pt,letterpaper]{article}
\usepackage[utf8]{inputenc}
\usepackage[spanish]{babel}
\usepackage{amsmath}
\usepackage{amsfonts}
\usepackage{amssymb}
\usepackage[margin=2.5cm]{geometry}

\begin{document}

\begin{center}
    \Large \textbf{Evaluación de Examen 2: Unidad II} \\
    \large Física para Ciencias de la Salud (005-1713) \\
    \vspace{0.5cm}
    \textbf{Estudiante:} Luisannys García \quad \textbf{C.I.:} 33.141.963
    \vspace{0.3cm} \\
    \large \textbf{Calificación Final: 8,5/10 puntos}
\end{center}

\vspace{0.5cm}

\section*{Análisis de Resultados}

\subsection*{1. Cinemática (Aceleración media)}
\textbf{Procedimiento y resultado:} Extrae los datos y efectúa el cálculo analítico a la perfección ($a = 4 \text{ m/s}^2$). No obstante, al llevar este valor a su columna de "Notación Científica" y "Respuesta", escribe erróneamente $4 \times 10 \text{ m/s}^2$. Matemáticamente, $4 \times 10$ equivale a $40$, lo cual altera el resultado clínico o físico reportado. \\
\textbf{Veredicto:} \textbf{Parcialmente correcto} (Se penaliza el error de notación matemática en la respuesta final) (1,5/2 pts).

\subsection*{2. Dinámica (Leyes de Newton)}
\textbf{Procedimiento y resultado:} Muestra un despeje correcto a partir de la Segunda Ley de Newton ($a = \frac{F}{m}$). Sustituye los valores dividiendo $150 \text{ N}$ entre $25 \text{ kg}$, logrando el cálculo exacto de $6 \text{ m/s}^2$. Al igual que en la pregunta anterior, reporta el resultado final redactado como $6 \times 10 \text{ m/s}^2$ (que equivaldría a 60). \\
\textbf{Veredicto:} \textbf{Parcialmente correcto} (Penalización por error de notación final) (1,5/2 pts).

\subsection*{3. Trabajo Mecánico}
\textbf{Procedimiento y resultado:} Plantea la fórmula de trabajo y efectúa la multiplicación ($400 \text{ N} \cdot 0,8 \text{ m}$), obteniendo de forma exacta $320 \text{ J}$. En este caso, la transformación a notación científica sí es matemáticamente impecable ($3,2 \times 10^2 \text{ J}$). \\
\textbf{Veredicto:} \textbf{Correcto} (2/2 pts).

\subsection*{4. Energía Potencial}
\textbf{Procedimiento y resultado:} Aplica a la perfección la fórmula de la energía potencial gravitatoria ($E_p = m \cdot g \cdot h$). El producto de los tres factores es matemáticamente exacto ($17,64 \text{ J}$) y su conversión a notación científica ($1,764 \times 10^1 \text{ J}$) es completamente válida y correcta. \\
\textbf{Veredicto:} \textbf{Correcto} (2/2 pts).

\subsection*{5. Potencia Mecánica}
\textbf{Procedimiento y resultado:} Plantea correctamente la fórmula de potencia mecánica ($P = \frac{W}{t}$) y efectúa la división aritmética de $75 / 60$ obteniendo el valor correcto de $1,25 \text{ W}$. Reincide en el error de notación al reportar su respuesta final como $1,25 \times 10 \text{ W}$ (que equivaldría a $12,5 \text{ W}$). \\
\textbf{Veredicto:} \textbf{Parcialmente correcto} (Penalización por error de notación en la conclusión) (1,5/2 pts).

\vspace{0.5cm}
\section*{Comentarios Generales y Sugerencias}
El examen evidencia un nivel de estructuración gráfica y procedimental sobresaliente. La creación de tablas separando rigurosamente cada fase del cálculo es una excelente iniciativa digna de elogio. Sus operaciones aritméticas son de hecho perfectas ($4$, $6$, $320$, $17,64$ y $1,25$).
\begin{itemize}
    \item Existe una confusión conceptual importante al usar la \textbf{Notación Científica}. Un número por sí solo no requiere que se le agregue "$\times 10$" al final si no estamos moviendo la coma decimal.
    \item Escribir $4 \times 10$ altera el número (lo convierte en $40$). Si se desea forzar un número de una sola cifra a tener base diez, el exponente debe ser estrictamente cero ($4 \times 10^0$), recordando que toda potencia elevada a la cero es 1 ($4 \times 1 = 4$).
    \item Se invita a la estudiante a revisar esta propiedad de los exponentes, manteniendo el mismo esquema brillante de organización en sus próximas asignaciones.
\end{itemize}

\end{document}